\section{Uniform exclusions and necessary source complexity}\label{sec:uniform}

The finite bound has constants independent of the number of selected
core families. This permits a numerical limit argument. It does not
permit an infinite source bank or a cut parameter growing with dimension.

\subsection{Original hop and variation profiles}

Fix $c,\beta>0$, $N\ge1$ and an integer $q\ge1$. In each admissible
geometry let $\mathcal K_0$ be an original $\mathrm Q_2$ core and let
$\mathcal E$ contain at most $q$ other complete families. Let $a_0$ bound
the hop count of every family outside $\mathcal K_0\cup\mathcal E$.
Sort the original core by nonincreasing hop count, with a fixed rule
for ties, and denote it by $G_1,G_2,\ldots$. Both profiles below are
zero at indices with no family. The variation profile uses this same ordering, not an ordering
by variation.

\begin{definition}[Near-full sequence and joint profile]\label{def:joint-profile}
A \emph{near-full sequence} consists of admissible fixed geometries and
nonempty correctness sets indexed by $\nu$, with common $c,\beta,N,q$,
such that
\begin{equation}\label{eq:near-full}
 d_\nu\longrightarrow\infty,\qquad
 a_{0,\nu}/L_\nu\longrightarrow0,\qquad
 \Delta_\nu/L_\nu\longrightarrow0.
\end{equation}
Its original profiles are
\[
 h_i^{(\nu)}=|G_i|/L_\nu,\qquad
 v_i^{(\nu)}=\Var(x_\nu;[p_{G_i},U_{G_i}])/L_\nu.
\]
A \emph{joint profile} is a coordinatewise subsequential limit
$(h_i,v_i)_{i\ge1}$ of both sequences of profiles. Write
\[
 \tau_R=\sum_{i>R}h_i,\qquad V_K=\sum_{i>K}v_i.
\]
\end{definition}

Such a further subsequence always exists: the family accounting in
\cref{sec:source-engine} gives
$\sum_i h_i^{(\nu)}\le c$ and $\sum_i v_i^{(\nu)}\le2c$.
Successive bounded-coordinate extractions and a diagonal subsequence
give simultaneous convergence. Passing each finite partial sum to the
limit gives $\sum_i h_i\le c$ and $\sum_i v_i\le2c$.
In particular both tails tend to zero. This does not assert that tails
of the finite profiles converge or that family topologies stabilize.

\begin{theorem}[Theorem U: necessary hop tails]\label{thm:necessary-hop-tail}
\notformalized
For every joint profile of every near-full sequence and all integers
$K\ge0$, $R\ge K$, $M\ge\max(2,R+1)$,
\begin{equation}\label{eq:necessary-hop-tail}
 \frac{\beta\tau_R}{M-R}
 \ge \frac{(1-V_K)_+^2}{4(K+1)(2M+4)^q}.
\end{equation}
Consequently, whenever $V_K<1$ and $R\ge\max(K,1)$,
\begin{equation}\label{eq:tail-floor}
 \tau_R\ge
 \frac{(1-V_K)^2}{\beta4^{q+1}(K+1)(R+2)^{q-1}}.
\end{equation}
If $q\ge2$, every such profile satisfies
$\sum_{i\ge1}i^{q-1}h_i=\infty$.
\end{theorem}
\begin{informalproof}
Fix a joint-profile subsequence and $K,R,M$. For each fixed integer
$J\ge R$, select the first $J$ original core families, distinguish its
first $K$, protect its first $R$, and rebuild the forest and ranks of
that finite selection. Absent families are simply omitted. Sorted hop
counts and their total bound imply that every omitted core family has
at most $cL_\nu/(J+1)$ hops. The finite theorem therefore applies with
\[
 a_\nu=\max\{a_{0,\nu},\lceil cL_\nu/(J+1)\rceil\},\qquad
 w_\nu=\lceil\lambda L_\nu\rceil,
\]
where $0<\lambda<1$ is also fixed before $\nu$ tends to infinity.
The finite-prefix quantities satisfy
\[
 \alpha_\nu\to V_{K,J}:=\sum_{K<i\le J}v_i,\qquad
 \tau_\nu/L_\nu\to\tau_{R,J}:=\sum_{R<i\le J}h_i.
\]
The normalized residual bill is bounded in the limit by
$\beta NC_qc^2/[\lambda(J+1)]$; the normalized movement, outer-lower
and stripe errors vanish, and $2(w_\nu-1)/L_\nu\to2\lambda$.
Using \cref{cor:orbit-count-deficit} and \cref{eq:near-full}, we obtain
\begin{equation}\label{eq:finite-prefix-limit}
 \frac{(1-V_{K,J})_+^2}{4(K+1)(2M+4)^q}
 \le\frac{\beta\tau_{R,J}}{M-R}+2\lambda
       +\frac{\beta NC_qc^2}{\lambda(J+1)}.
\end{equation}
This is an inequality for every finite $J$. Now let $J$ tend to infinity
in these numerical inequalities, and then let $\lambda$ decrease to zero.
The constant $C_q$ depends on $M,q$, not $J$, so this proves
\cref{eq:necessary-hop-tail}.

Taking $M=2R+2$ gives $2M+4=4(R+2)$ and proves
\cref{eq:tail-floor}. A finite $K$ with $V_K\le1/2$ always exists.
If $\sum_i i^{q-1}h_i$ were finite, then for $R\ge1$
\[
 (R+2)^{q-1}\tau_R
 \le2^{q-1}\sum_{i>R}i^{q-1}h_i\longrightarrow0,
\]
contrary to \cref{eq:tail-floor} for that $K$.
Every step used a finite certificate. Omitted families remained in the
residual bill until after the dimension limit; no family was discarded
by an exchange of limits.
\end{informalproof}

\begin{corollary}[One arbitrary exceptional family]\label{cor:one-exception}
\notformalized
Fix $c,\beta>0$ and $N\ge1$. There exist $\eps>0$ and $d_0$ such that
every admissible geometry with a $\mathrm Q_2$ core, at most one other
complete exceptional family, $d\ge d_0$, and $a_0\le\eps L$ satisfies
\begin{equation}\label{eq:uniform-deficit}
 \Delta\ge\eps L,\qquad
 |\mathcal C|\le2^{(1-\eps)L}.
\end{equation}
No envelope on the sorted core-family sizes is assumed.
\end{corollary}
\begin{informalproof}
If no such constants existed, choose a violating geometry for each
$\eps=1/n$ and $d_0=n$. These form a near-full sequence and have a
joint profile. Choose $K$ with $V_K\le1/2$. For $q=1$,
\cref{eq:tail-floor} bounds $\tau_R$ below by a fixed positive constant
for every $R\ge\max(K,1)$, contrary to $\tau_R\to0$.
\end{informalproof}

\begin{corollary}[Common size envelopes]\label{cor:common-envelope}
\notformalized
Fix $c,\beta,N,q$ as above and a nonnegative summable sequence $(g_i)$
such that
\begin{equation}\label{eq:envelope-test}
 \liminf_{R\to\infty}(R+2)^{q-1}\sum_{i>R}g_i=0.
\end{equation}
There exist $\eps>0,d_0$ for which \cref{eq:uniform-deficit} holds
whenever $d\ge d_0$, $a_0\le\eps L$ and $|G_i|\le Lg_i$ for all
original core families in their sorted order. In particular every
common power envelope $g_i=H/i^p$, with fixed $H>0$ and $p>q$, is covered.
\end{corollary}
\begin{informalproof}
Negating the uniform conclusion again gives a near-full sequence.
In a joint profile $h_i\le g_i$ for each $i$, so
\cref{eq:envelope-test} contradicts \cref{eq:tail-floor} after choosing
$K$ with $V_K<1$. For a power envelope,
$\sum_{i>R}H/i^p\le H R^{1-p}/(p-1)$; multiplication by
$(R+2)^{q-1}$ tends to zero when $p>q$.
\end{informalproof}

\begin{corollary}[A fixed finite skeleton]\label{cor:finite-skeleton}
\notformalized
For fixed $c,\beta,N$ and $q\ge1$, there are $\eps>0,d_0$ such that
\cref{eq:uniform-deficit} holds for every admissible geometry in which
all but at most $q$ complete families have at most $a_0$ hops,
$a_0\le\eps L$, and $d\ge d_0$.
\end{corollary}
\begin{informalproof}
Take the core empty and use $g_i=0$ in \cref{cor:common-envelope}.
The full geometry and its outside-family bound are unchanged.
\end{informalproof}

These are restrictions on a common full geometry. In the one-exception
case they settle the paid-deficit comparison without a coding conjecture.
For $q=2$, the necessary tail floor is of order $1/R$, compatible with
an inverse-square size profile; a further geometric argument is needed.
The proof supplies uniform constants by contradiction, not effective
numerical values of $\eps$ and $d_0$.

\subsection{Necessary fractional-cover complexity}

Along a fixed joint-profile subsequence, select the first $J$ families
and use $M=2R+2$ as above. For each phase and column put
$\kappa_{\nu,J,R,\rho,j}=\kappa_\rho(S_\rho(s_j))$.
All these quantities are finite and at least one. Write
\begin{equation}\label{eq:reciprocal-cover-limit}
 Q_R^*=\limsup_{J\to\infty}\limsup_{\nu\to\infty}
                \mathbb E_{\rho,j}\frac1{\kappa_{\nu,J,R,\rho,j}},
\end{equation}
where the expectation is the arithmetic mean over legal phases and
columns. This order of limits is part of the definition.

\begin{theorem}[Necessary reciprocal-cover bound]\label{thm:necessary-cover}
\notformalized
For every near-full sequence, every joint profile, any $K$ with $V_K<1$
and every $R\ge\max(K,1)$,
\begin{equation}\label{eq:necessary-cover}
 Q_R^*\le
 \frac{4\beta(K+1)\tau_R}{(1-V_K)^2(R+2)}.
\end{equation}
This holds for any fixed number $q$ of exceptions.
If $|G_i|\le HL/i^2$ and $V_K\le1/2$, then for every fixed $s\ge1$,
\begin{equation}\label{eq:cover-saturation}
 \limsup_J\limsup_\nu
 \Pr_{\rho,j}\{\kappa_{\nu,J,R,\rho,j}\le s\}
 \le \frac{16\beta(K+1)Hs}{R(R+2)}.
\end{equation}
\end{theorem}
\begin{informalproof}
Use the exact witnesses in \cref{thm:paid-source-cover}. For fixed
$J,R,K,\lambda$ the geometric coefficient tends to
$(1-V_{K,J})^2/[4(K+1)]>0$. The same error calculation as in
\cref{eq:finite-prefix-limit} bounds the limsup of the mean reciprocal
cover by that coefficient's reciprocal times
\[
 \frac{\beta\tau_{R,J}}{R+2}+2\lambda
       +\frac{\beta NC_qc^2}{\lambda(J+1)}.
\]
Take the limsup in $J$, then let $\lambda$ decrease to zero. This proves
\cref{eq:necessary-cover}. The elementary pointwise inequality
$\mathbf1_{\{\kappa\le s\}}\le s/\kappa$ and
$\tau_R\le H/R$ give \cref{eq:cover-saturation}.
\end{informalproof}

\begin{corollary}[Uniform cover criteria]\label{cor:uniform-cover}
\notformalized
Fix $c,\beta,N,q$. Each of the following defines a class with constants
$\eps>0,d_0$ for which \cref{eq:uniform-deficit} holds when
$d\ge d_0$ and $a_0\le\eps L$.
\begin{enumerate}[label=(\alph*)]
\item For some common constant $0<B<\infty$, selecting the first $J$
original core families and protecting the first $R$, for every
$J\ge R\ge1$ and $M=2R+2$, gives
$\mathbb E_{\rho,j}\kappa_{J,R,\rho,j}\le B(R+2)$.
No size envelope is required.
\item For fixed $H>0,p>1$ the original sizes satisfy $|G_i|\le HL/i^p$,
and for a common positive function $f(R)=o(R^p)$ every such finite
prefix selection satisfies $\mathbb E_{\rho,j}\kappa_{J,R,\rho,j}\le f(R)$.
\end{enumerate}
Here $\kappa_{J,R,\rho,j}$ is the fractional cover of the selected
geometry's complete source set at the first output time of column $j$.
The bounds are uniform over the class and dimension; $\eps,d_0$ may
depend on all the fixed class data. In (b), the
inverse-square case requires mean cover $o(R^2)$.
\end{corollary}
\begin{informalproof}
Suppose either uniform conclusion fails and extract a joint profile of
a near-full sequence. Jensen's inequality for $u\mapsto1/u$ gives,
respectively, $Q_R^*\ge1/[B(R+2)]$ or $Q_R^*\ge1/f(R)$.
Choose $K$ with $V_K<1$. In (a), \cref{eq:necessary-cover} then forces
$\tau_R$ to stay above a positive constant, a contradiction.
In (b), $\tau_R\le HR^{1-p}/(p-1)$ bounds its right side by
a fixed multiple of $R^{-p}$, contradicting $f(R)=o(R^p)$.
\end{informalproof}

The saturation bound also rules out a misleading average picture:
in a near-full inverse-square sequence, rare very complicated columns
cannot compensate for a positive fraction with cover complexity
$o(R^2)$, in the stated ordered limits. This is a necessary condition,
not a construction of an economical cover.

\subsection{The movement of complete older observables}

For $q=2$, use one cluster for all nonroot elements of $F_\rho(s_j)$
in \cref{prop:continuation-covers}(c), and use $E=0$ when that set is
empty or a singleton. Let $e_{\nu,J,R}$ be the mean of $E/L_\nu$
over phases and columns, with $M=2R+2$, and put
$e_R^*=\liminf_J\liminf_\nu e_{\nu,J,R}$.

\begin{corollary}[Necessary complete-observable movement]\label{cor:necessary-exposure}
\notformalized
Along every subsequence realizing a joint profile of a two-exception
near-full sequence, for every integer $K\ge0$ and $R\ge\max(K,1)$,
\begin{equation}\label{eq:necessary-exposure}
 e_R^*\ge
 \frac{(1-V_K)_+^2}{32(K+1)(R+2)}-\frac{\beta\tau_R}{R+2}.
\end{equation}
For fixed $K$ with $V_K<1$ this implies
$\liminf_{R\to\infty}(R+2)e_R^*\ge(1-V_K)^2/[32(K+1)]$.
\end{corollary}
\begin{informalproof}
The one-cluster estimate, also valid when there are no nonroot seeds,
is $(1-\alpha)_+^2/[8(K+1)(2M+4)]-E/L$.
Insert it in \cref{thm:paid-source-cover} and move the mean $E/L$ to
the other side. Take the liminf in $\nu$, then in $J$, and then let
$\lambda$ decrease to zero exactly as in
\cref{eq:finite-prefix-limit}. Substitution of $M=2R+2$ gives
\cref{eq:necessary-exposure}; $\tau_R\to0$ gives the last assertion.
\end{informalproof}

Here the moving observables include every older continuation, not just
the immediate input coordinate at an exceptional landing. This is the
quantity that a successful approximation argument would have to control.
